\documentclass[border=10pt]{standalone}
\usepackage{tikz,ctex}
\usetikzlibrary{positioning}

\begin{document}

\begin{tikzpicture}[
    % 全局设置
    node distance = 1.6cm and 1.4cm,
    every node/.style = {
        draw,
        rounded corners,
        thick,
        align = center,
        font = \sffamily\small,
        minimum width = 2.0cm,
        minimum height = 0.9cm
    },
    % 样式定义
    centerStyle/.style = {
        fill = blue!45,
        text = white,
        font = \sffamily\bfseries\large,
        line width = 2pt,
        minimum size = 2.5cm,
        circle
    },
    branchTitle/.style = {
        font = \sffamily\bfseries,
        line width = 1.5pt
    },
    subNode/.style = {
        fill = white,
        font = \sffamily\small,
        minimum width = 2.0cm,
        minimum height = 0.8cm
    },
    leafNode/.style = {
        fill = white,
        font = \sffamily\footnotesize,
        minimum width = 1.8cm,
        minimum height = 0.7cm
    },
    arrowStyle/.style = {
        ->,
        thick,
        >=stealth,
        shorten >=2pt
    }
]

    % ================= 1. 中心节点 =================
    \node[centerStyle] (center) {7.3.收敛性 \\ Convergence};

    % ================= 2. 上方分支：学习率与曲率 (红) =================
    \node[branchTitle, fill=red!15, above=of center] (lr) {学习率与曲率 \\ Learning Rate \& Curvature};

    \node[subNode, above=of lr, xshift=-6cm] (valley) {狭长山谷 \\ Narrow Valley};
    \node[subNode, above=of lr, xshift=-2cm] (oscillate) {振荡发散 \\ Oscillation};
    \node[subNode, above=of lr, xshift=2cm] (eigen) {特征值分解 \\ Eigenvalues};
    \node[subNode, above=of lr, xshift=6cm] (condition) {条件数 \\ Condition Number};

    % ================= 3. 右侧分支：线性收敛分析 (蓝) =================
    \node[branchTitle, fill=orange!15, right=of center] (linear) {线性收敛分析 \\ Linear Convergence};

    \node[subNode, right=of linear, yshift=3cm] (decouple) {方向解耦 \\ Decoupled Directions};
    \node[subNode, right=of linear, yshift=1cm] (factor) {收敛因子 \\ Convergence Factor};
    \node[subNode, right=of linear, yshift=-1cm] (bound) {上界 $\eta < 2/\lambda_{\max}$ \\ Upper Bound};
    \node[subNode, right=of linear, yshift=-3cm] (slow) {最慢方向 \\ Slowest Direction};

    % ================= 4. 下方分支：动量法 (绿) =================
    \node[branchTitle, fill=green!15, below=of center] (momentum) {动量法 \\ Momentum};

    \node[subNode, below=of momentum, xshift=-6cm] (inertia) {惯性项 \\ Inertia};
    \node[subNode, below=of momentum, xshift=-3cm] (smooth) {平滑振荡 \\ Smooth Oscillations};
    \node[subNode, below=of momentum, xshift=0cm] (effective) {有效学习率 \\ Effective LR};
    \node[subNode, below=of momentum, xshift=3cm] (mu) {动量参数 $\mu$ \\ Momentum Parameter};
    \node[subNode, below=of momentum, xshift=7cm] (nesterov) {Nesterov 动量 \\ Nesterov Momentum};

    % ================= 5. 左侧分支：自适应学习率 (紫) =================
    \node[branchTitle, fill=purple!15, left=of center] (adaptive) {自适应学习率 \\ Adaptive Learning Rates};

    \node[subNode, left=of adaptive, yshift=4cm] (schedule) {学习率调度 \\ LR Schedule};
    \node[subNode, left=of adaptive, yshift=2cm] (adagrad) {AdaGrad \\ 累积平方梯度};
    \node[subNode, left=of adaptive, yshift=0cm] (rmsprop) {RMSProp \\ 指数加权平均};
    \node[subNode, left=of adaptive, yshift=-2cm] (adam) {Adam \\ 自适应矩估计};
    \node[subNode, left=of adaptive, yshift=-4cm] (bias) {偏差校正 \\ Bias Correction};

    % ================= 连线逻辑 =================
    % 中心 -> 四大分支标题
    \draw[arrowStyle, red] (center) -- (lr);
    \draw[arrowStyle, blue] (center) -- (linear);
    \draw[arrowStyle, green!60!black] (center) -- (momentum);
    \draw[arrowStyle, purple] (center) -- (adaptive);

    % 分支标题 -> 子节点 (上方)
    \draw[arrowStyle, red!60] (lr) -- (valley.south);
    \draw[arrowStyle, red!60] (lr) -- (oscillate.south);
    \draw[arrowStyle, red!60] (lr) -- (eigen.south);
    \draw[arrowStyle, red!60] (lr) -- (condition.south);

    % 分支标题 -> 子节点 (右侧)
    \draw[arrowStyle, blue!60] (linear) -- (decouple.west);
    \draw[arrowStyle, blue!60] (linear) -- (factor.west);
    \draw[arrowStyle, blue!60] (linear) -- (bound.west);
    \draw[arrowStyle, blue!60] (linear) -- (slow.west);

    % 分支标题 -> 子节点 (下方)
    \draw[arrowStyle, green!60!black] (momentum) -- (inertia.north);
    \draw[arrowStyle, green!60!black] (momentum) -- (smooth.north);
    \draw[arrowStyle, green!60!black] (momentum) -- (effective.north);
    \draw[arrowStyle, green!60!black] (momentum) -- (mu.north);
    \draw[arrowStyle, green!60!black] (momentum) -- (nesterov.north);

    % 分支标题 -> 子节点 (左侧)
    \draw[arrowStyle, purple!60] (adaptive) -- (schedule.east);
    \draw[arrowStyle, purple!60] (adaptive) -- (adagrad.east);
    \draw[arrowStyle, purple!60] (adaptive) -- (rmsprop.east);
    \draw[arrowStyle, purple!60] (adaptive) -- (adam.east);
    \draw[arrowStyle, purple!60] (adaptive) -- (bias);

\end{tikzpicture}

\end{document}